Kurs:Algebraische Kurven (Osnabrück 2012)/Vorlesung 29/latex
Kurs:Algebraische Kurven (Osnabrück 2012)/Vorlesung 29/latex

\setcounter{section}{29}


  \zwischenueberschrift{Projektion weg von einem Punkt}


\inputdefinition

{}

{


Die
\definitionsverweis {Abbildung}{}{}
\maabbeledisp {} { {\mathbb P}^{ n }_{ K} \setminus \{ (1,0 , \ldots , 0) \}
} { {\mathbb P}^{n-1}_{K}
} { ( x_0 , x_1, \ldots ,  x_n)
} { (x_1 , \ldots , x_n)
} {,}
heißt \definitionswort {die Projektion weg vom Punkt}{}

\mathl{(1,0 , \ldots , 0)}{.}


}


Diese Abbildung ist ein wohldefinierter Morphismus, der außerhalb des \stichwort {Zentrums} {}

\mathl{(1,0 , \ldots , 0)}{} der Projektion definiert ist. Jedem anderen Punkt wird derjenige Punkt des

\mathl{{\mathbb P}^{n-1}_{K}}{} zugeordnet, der der Geraden durch den Punkt und dem Zentrum entspricht. Daher ist die Abbildung surjektiv und jede Faser ist eine projektive Gerade ohne den Zentrumspunkt, also eine affine Gerade
\zusatzklammer {es liegt ein sogenanntes \stichwort {Geradenbündel} {} über dem  ${\mathbb P}^{n-1}_{K}$ vor} {} {.}
Es handelt sich um die Fortsetzung der Kegelabbildung
\maabbdisp {} {{ {\mathbb A}_{ K }^{ n  } } \setminus \{0\}} { {\mathbb P}^{n-1}_{K}
} {}
auf den punktierten projektiven Raum. Die entsprechende Abbildung kann man zu jedem Zentrumspunkt definieren, siehe
Aufgabe 29.7.


  \zwischenueberschrift{Abbildungen nach ${\mathbb P}^{1}_{K}$}


Der folgende Satz liefert eine neue Version
der Noetherschen Normalisierung.


\inputfaktbeweis

{Ebene projektive Kurve/Abbildung nach P^1 über Projektion von einem Punkt/Fakt}

{Satz}

{}

{


\faktsituation {}

\faktvoraussetzung {Es sei $K$ ein
\definitionsverweis {algebraisch abgeschlossener Körper}{}{}
und sei


\mavergleichskette

{\vergleichskette

{C
}

{ \subset }{ {\mathbb P}^{2}_{K}
}

{  }{
}

{    }{
}

{   }{
}

}
{}{}{}
eine
\definitionsverweis {ebene projektive Kurve}{}{}
vom Grad $d$.}

\faktfolgerung {Dann gibt es einen surjektiven
\definitionsverweis {Morphismus}{}{}
\maabbdisp {} { C } { {\mathbb P}^{1}_{K}
} {}
derart, dass alle
\definitionsverweis {Fasern}{}{}
aus maximal $d$ Punkten bestehen.}

\faktzusatz {}

\faktzusatz {}


}

{


Es sei


\mavergleichskette

{\vergleichskette

{ P
}

{ \in }{ {\mathbb P}^{2}_{K}
}

{  }{
}

{    }{
}

{   }{
}

}
{}{}{}
ein Punkt, der nicht auf der Kurve liegt. Einen solchen Punkt gibt es, da der Körper insbesondere unendlich ist. Wir betrachten die
\definitionsverweis {Projektion weg von $P$}{}{,}
die insgesamt einen Morphismus


\mathdisp {C \hookrightarrow {\mathbb P}^{2}_{K} \setminus \{P\} \longrightarrow {\mathbb P}^{1}_{K}} {  }


induziert. Die Faser dieses Morphismus über einem Punkt


\mavergleichskette

{\vergleichskette

{ Q
}

{ \in }{ {\mathbb P}^{1}_{K}
}

{  }{
}

{    }{
}

{   }{
}

}
{}{}{}
\zusatzklammer {der eine Richtung in


\mavergleichskettek

{\vergleichskettek

{ P
}

{ \in }{ {\mathbb P}^{2}_{K}
}

{  }{
}

{    }{
}

{   }{
}

}
{}{}{}
repräsentiert} {} {}
besteht genau aus den Punkten der Kurve, die auf der durch $Q$ definierten Geraden


\mavergleichskettedisp

{\vergleichskette

{ G
}

{ =} { V_+(aX+bY+cZ)
}

{ \cong} { {\mathbb P}^{1}_{K}
}

{ \subset} { {\mathbb P}^{2}_{K}
}

{ } {
}

}
{}{}{}
liegen. Daher wird die Faser über $Q$ auf $G$ beschrieben, indem man in der Kurvengleichung


\mavergleichskette

{\vergleichskette

{ C
}

{ = }{ V_+(F)
}

{  }{
}

{    }{
}

{   }{
}

}
{}{}{}
mittels der Geradengleichung eine Variable eliminiert. Das Ergebnis ist ein homogenes Polynom $\bar{F}$ in zwei Variablen vom Grad $d$, das nicht $0$ ist, denn sonst wäre $P$ ein Punkt der Kurve. Da wir über einem algebraisch abgeschlossenen Körper sind, besitzt dieses Polynom $\bar{F}$ mindestens eine und höchstens $d$ Nullstellen, die alle von $P$ verschieden sind. Dies ergibt die Surjektivität und die Abschätzung für die Faser.


}


\inputfaktbeweis

{Glatte projektive Kurven/Rationale Funktion als Morphismus nach P^1/Fakt}

{Satz}

{}

{


\faktsituation {}

\faktvoraussetzung {Es sei $K$ ein
\definitionsverweis {Körper}{}{}
und


\mavergleichskette

{\vergleichskette

{C
}

{ \subset }{ {\mathbb P}^{2}_{K}
}

{  }{
}

{    }{
}

{   }{
}

}
{}{}{}
eine
\definitionsverweis {glatte}{}{}
irreduzible
\definitionsverweis {ebene projektive Kurve}{}{.}
Es sei


\mavergleichskette

{\vergleichskette

{D
}

{ = }{ C \cap D_+(Z)
}

{ \cong }{ K\!-\!\operatorname{Spek}\, \left(  R  \right)
}

{    }{
}

{   }{
}

}
{}{}{}
ein affines Teilstück davon. Es sei


\mavergleichskette

{\vergleichskette

{ q
}

{ = }{ g/h
}

{ \in }{ Q(R)
}

{    }{
}

{   }{
}

}
{}{}{}
eine rationale Funktion
\zusatzklammer {mit \mathlk{g,h \in R,\, h \neq 0}{}} {} {.}}

\faktfolgerung {Dann gibt es einen eindeutigen
\definitionsverweis {Morphismus}{}{}
\maabbdisp {\varphi} { C } { {\mathbb P}^{1}_{K}
} {}
derart, dass das Diagramm


\mathdisp {\begin{matrix} D(h) & \stackrel{ g/h }{\longrightarrow} &  {\mathbb A}^{1}_{K} \cong D_+(s) &  \\ \downarrow & & \downarrow  & \\  C  & \stackrel{ \varphi }{\longrightarrow} &  {\mathbb P}^{1}_{ K }  & \! \! \! \!\!\!\!\!   \\ \end{matrix}} {  }


kommutiert.}

\faktzusatz {}

\faktzusatz {}


}

{


Wir definieren zunächst auf $D$ eine Fortsetzung
\maabbdisp {\varphi} { D } { {\mathbb P}^{1}_{}
} {}
der rationalen Funktion

\mathl{g/h}{.} Es sei hierzu


\mavergleichskette

{\vergleichskette

{ P
}

{ \in }{ D
}

{  }{
}

{    }{
}

{   }{
}

}
{}{}{}
ein Punkt der Kurve. Bei


\mavergleichskette

{\vergleichskette

{ P
}

{ \in }{ D(h)
}

{  }{
}

{    }{
}

{   }{
}

}
{}{}{}
ist nichts zu tun, sei also


\mavergleichskette

{\vergleichskette

{ h(P)
}

{ = }{ 0
}

{  }{
}

{    }{
}

{   }{
}

}
{}{}{.}
Da die Kurve glatt ist, ist nach
Satz 23.6
der lokale Ring $B$ der Kurve im Punkt $P$ ein
\definitionsverweis {diskreter Bewertungsring}{}{.}
Daher hat der Quotient

\mathl{g/h}{} dort eine Beschreibung als


\mavergleichskettedisp

{\vergleichskette

{ \frac{g}{h}
}

{ =} { u \pi^n
}

{ } {
}

{ } {
}

{ } {
}

}
{}{}{}
mit \mathkon  { u \in B^\times  }  {  und  }  { n \in \Z  }{ }
\zusatzklammer {$\pi$ sei eine Ortsuniformisierende} {} {.}
Es gibt eine offene Umgebung


\mavergleichskette

{\vergleichskette

{P
}

{ \in }{ D(\psi )
}

{ \subseteq }{ D
}

{    }{
}

{   }{
}

}
{}{}{}
derart, dass $\pi$ und $u$ über

\mathl{D( \psi)}{} definiert sind und $u$ dort eine Einheit ist. Bei


\mavergleichskette

{\vergleichskette

{n
}

{ \geq }{ 0
}

{  }{
}

{    }{
}

{   }{
}

}
{}{}{}
ist


\mavergleichskette

{\vergleichskette

{ g/h
}

{ \in }{ R_\psi
}

{  }{
}

{    }{
}

{   }{
}

}
{}{}{}
und die Undefiniertheitsstelle ist also sogar als Abbildung nach

\mathl{{\mathbb A}^{1}_{K}}{} \anfuehrung{hebbar}{.} Bei


\mavergleichskette

{\vergleichskette

{ n
}

{ \leq }{ 0
}

{  }{
}

{    }{
}

{   }{
}

}
{}{}{}
ist der umgekehrte Bruch


\mavergleichskette

{\vergleichskette

{ h/g
}

{ = }{ u^{-1} \pi^{-n}
}

{  }{
}

{    }{
}

{   }{
}

}
{}{}{}
auf

\mathl{D(\psi)}{} definiert als eine Abbildung nach

\mathl{{\mathbb A}^{1}_{K}}{.} Mittels der \anfuehrung{verdrehten Einbettung}{}

\mathl{{\mathbb A}^{1}_{K} \cong D_+(t) \hookrightarrow {\mathbb P}^{1}_{K}}{} erhält man eine Abbildung nach

\mathl{{\mathbb P}^{1}_{K}}{.}


Wir müssen zeigen, dass diese zwei Morphismen in die projektive Gerade dort, wo beide definiert sind, übereinstimmen. Das sind die Punkte $P$, in denen $\varphi$ weder eine Nullstelle noch einen Pol besitzt. Die Verträglichkeit folgt daraus, dass in einer offenen Umgebung


\mavergleichskette

{\vergleichskette

{ P
}

{ \in }{ U
}

{  }{
}

{    }{
}

{   }{
}

}
{}{}{}
eine Abbildung
\maabbdisp {g/h} { U } { ({\mathbb A}^{1}_{K})^\times = {\mathbb A}^{1}_{K} \setminus \{0\}
} {}
vorliegt, und dass das Diagramm


\mathdisp {\begin{matrix}  ({\mathbb A}^{1}_{K}) ^\times  & \stackrel{ i^{-1} }{\longrightarrow} &  {\mathbb A}^{1}_{K} \cong D_+(t) &  \\ \downarrow & & \downarrow  & \\  {\mathbb A}^{1}_{K} \cong D_+(s)    & \stackrel{  }{\longrightarrow} &  {\mathbb P}^{1}_{ K }  & \! \! \! \!\!\!\!\!   \\ \end{matrix}} {  }


kommutiert. Dies ergibt einen wohldefinierten Morphismus auf dem affinen Stück $D$.


Für einen beliebigen Punkt der projektiven Kurve $C$ und eine affine Umgebung


\mavergleichskette

{\vergleichskette

{ P
}

{ \in }{ D'
}

{ \subset }{ C
}

{    }{
}

{   }{
}

}
{}{}{}
liegt die gleiche Situation vor, da


\mavergleichskette

{\vergleichskette

{ D_+(h) \cap D'
}

{ \neq }{ \emptyset
}

{  }{
}

{    }{
}

{   }{
}

}
{}{}{}
ist, und somit auf einer offenen nichtleeren Menge die rationale Funktion
\zusatzklammer {mit anderen Zählern und Nennern} {} {}
definiert ist. Damit lässt sich das vorhergehende Argument genauso anwenden.


Die Eindeutigkeit folgt daraus, dass auf jeder affinen offenen Menge


\mavergleichskette

{\vergleichskette

{ P
}

{ \in }{ U
}

{  }{
}

{    }{
}

{   }{
}

}
{}{}{}
der Durchschnitt

\mathl{U \cap D_+(h)}{} nicht leer ist. Ein Morphismus auf einer integren Varietät in die affine Gerade

\mathl{{\mathbb A}^{1}_{K}}{} ist durch die rationale Funktion eindeutig festgelegt.


}


\inputbeispiel{}

{


Die Inversenbildung
\maabbeledisp {} { {\mathbb A}^{1}_{K} \supset D(z)} { {\mathbb A}^{1}_{K} } { z } {z^{-1} } {,}
lässt sich zu einem bijektiven Morphismus
\maabbeledisp {} { {\mathbb P}^{1}_{K} } { {\mathbb P}^{1}_{K} } {(x,y)} {(y,x) } {,}
fortsetzen. Dies folgt direkt aus
Satz 29.3.
Dabei geht ein Punkt


\mavergleichskette

{\vergleichskette

{ z
}

{ \neq }{ 0
}

{  }{
}

{    }{
}

{   }{
}

}
{}{}{}
auf

\mathl{1/z}{} und der Nullpunkt geht auf den unendlich fernen Punkt $\infty$.


}


  \zwischenueberschrift{Parametrisierte projektive ebene Kurven}


Sei eine rational parametrisierte Kurve

\mathl{s \mapsto \left( \frac{\varphi_1(s)}{\psi(s)}, \frac{\varphi_2(s)}{ \psi(s)} \right)}{} gegeben. Wir haben in
Satz 6.11
gesehen, dass das Bild eine algebraische Gleichung erfüllt. Dabei haben wir im dortigen Beweis schon die homogenisierte Parametrisierung verwendet, die jetzt als projektive Fortsetzung wieder auftaucht.


\inputfaktbeweis

{Rationale Kurvenparametrisierung/Fortsetzung auf projektive Gerade/Fakt}

{Satz}

{}

{


\faktsituation {}

\faktvoraussetzung {Es sei
\maabbeledisp {} { {\mathbb A}^{1}_{K} \supseteq D(\psi) } { {\mathbb A}^{2}_{K} } { s } { \left( \frac{\varphi_1(s)}{ \psi(s)}, \frac{\varphi_2(s)}{ \psi(s)} \right)
} {,}
eine rationale Parametrisierung in gekürzter
\zusatzklammer {d.h. die \mathlk{\varphi_1, \varphi_2, \psi}{} haben keinen gemeinsamen Teiler} {} {}
Darstellung. Es sei $d$ der maximale Grad der beteiligten Polynome und es seien

\mathl{\hat{ \varphi_1   }, \hat{ \varphi_2   }, \hat{  \psi  }}{} die
\definitionsverweis {Homogenisierungen}{}{}
\zusatzklammer {bezüglich der neuen Variablen $t$} {} {}
davon. Es seien $H_i$ die Produkte dieser Homogenisierungen mit einer Potenz von $t$ derart, dass

\mathl{H_1,H_2,H_3}{} alle den Grad $d$ besitzen.}

\faktfolgerung {Dann definieren die

\mathl{H_1,H_2,H_3}{} einen Morphismus
\maabbeledisp {H} { {\mathbb P}^{1}_{K} } { {\mathbb P}^{2}_{K} } {(s,t)} { \left( H_1(s,t)  , \,  H_2(s,t)  , \,  H_3(s,t)                 \right)
} {,}
derart, dass das Diagramm


\mathdisp {\begin{matrix}  {\mathbb A}^{1}_{K} \supseteq D(\psi) & \stackrel{  }{\longrightarrow} &  {\mathbb A}^{2}_{K} \cong D_+(Z) &  \\ \downarrow & & \downarrow  & \\  {\mathbb P}^{1}_{K}  & \stackrel{ H }{\longrightarrow} &  {\mathbb P}^{2}_{K}  & \! \! \! \!\!\!\!\!   \\ \end{matrix}} {  }


kommutativ ist.}

\faktzusatz {Dabei liegt das Bild unter $H$ auf dem projektiven Abschluss der affinen Bildkurve.}

\faktzusatz {}


}

{


Die Abbildung $H$ ist aufgrund von
Aufgabe 29.6
wohldefiniert, und zwar auf ganz

\mathl{{\mathbb P}^{1}_{K}}{,} da

\mathl{\varphi_1, \varphi_2, \psi}{} insgesamt teilerfremd sind. Zur Kommutativität muss man lediglich beachten, dass


\mavergleichskette

{\vergleichskette

{s
}

{ \in }{ D(\psi)
}

{ \subseteq }{ {\mathbb A}^{1}_{K}
}

{    }{
}

{   }{
}

}
{}{}{}
einerseits über

\mathl{(s,1)}{} auf


\mavergleichskettedisp

{\vergleichskette

{ (H_1(s,1), H_2(s,1),H_3(s,1))
}

{ =} { \left(  \varphi_1(s)  , \,   \varphi_2(s)  , \,   \psi(s)                 \right)
}

{ } {
}

{ } {
}

{ } {
}

}
{}{}{}
abgebildet wird und andererseits auf


\mavergleichskettedisp

{\vergleichskette

{ \left( \frac { \varphi_1(s)}{ \psi(s)}, \frac{\varphi_2(s)}{ \psi(s)}  ,1 \right)
}

{ =} { \left(  \varphi_1(s)  , \,   \varphi_2(s)  , \,   \psi(s)                 \right)
}

{ } {
}

{ } {
}

{ } {
}

}
{}{}{.}
Für den Zusatz sei $C$ der affine Abschluss des Bildes und


\mavergleichskette

{\vergleichskette

{\bar{C}
}

{ \subset }{ {\mathbb P}^{2}_{K}
}

{  }{
}

{    }{
}

{   }{
}

}
{}{}{}
der
\definitionsverweis {projektive Abschluss}{}{}
davon. Wir betrachten das offene Komplement


\mavergleichskette

{\vergleichskette

{ U
}

{ = }{ {\mathbb P}^{2}_{K} \setminus \bar{C}
}

{  }{
}

{    }{
}

{   }{
}

}
{}{}{.}
Da die Abbildung stetig ist, ist das Urbild

\mathl{H^{-1}(U)}{} offen in

\mathl{{\mathbb P}^{1}_{K}}{,} und es kann nur Punkte aus

\mathl{{\mathbb P}^{1}_{K} \setminus D(\psi)}{} enthalten. Eine endliche und offene Teilmenge der projektiven Geraden muss aber leer sein.


}


\inputfaktbeweis

{Ebene projektive Kurve/Graph eines Polynoms in einer Variable/Singularität im Unendlichen/Fakt}

{Satz}

{}

{


\faktsituation {}

\faktvoraussetzung {Es sei $K$ ein
\definitionsverweis {algebraisch abgeschlossener Körper}{}{}
und sei


\mavergleichskette

{\vergleichskette

{ F
}

{ \in }{ K[X]
}

{  }{
}

{    }{
}

{   }{
}

}
{}{}{}
ein Polynom in einer Variablen vom Grad


\mavergleichskette

{\vergleichskette

{ d
}

{ \geq }{ 1
}

{  }{
}

{    }{
}

{   }{
}

}
{}{}{.}}

\faktfolgerung {Dann wird der
\definitionsverweis {projektive Abschluss}{}{}
$C$ des
\definitionsverweis {Graphen}{}{}


\mathl{V(Y-F(X))}{} durch

\mathl{V \left(  YZ^{d-1} - \hat{  F   }(X,Z)  \right)}{} beschrieben, wobei

\mathl{\hat{  F   } (X,Z)}{} die
\definitionsverweis {Homogenisierung}{}{}
von $F$ bezeichnet. Dabei gibt es in $C$ bei


\mavergleichskette

{\vergleichskette

{d
}

{ = }{ 1
}

{  }{
}

{    }{
}

{   }{
}

}
{}{}{}
\zusatzklammer {mit


\mavergleichskettek

{\vergleichskettek

{ F
}

{ = }{ aX+b
}

{  }{
}

{    }{
}

{   }{
}

}
{}{}{}} {} {}
noch den
\definitionsverweis {glatten Punkt}{}{}


\mathl{(1,a ,0)}{} und bei


\mavergleichskette

{\vergleichskette

{ d
}

{ \geq }{ 2
}

{  }{
}

{    }{
}

{   }{
}

}
{}{}{}
noch den Punkt

\mathl{(0,1,0)}{,} der bei


\mavergleichskette

{\vergleichskette

{ d
}

{ \geq }{ 3
}

{  }{
}

{    }{
}

{   }{
}

}
{}{}{}
singulär ist.}

\faktzusatz {Bei


\mavergleichskette

{\vergleichskette

{d
}

{ \geq }{ 2
}

{  }{
}

{    }{
}

{   }{
}

}
{}{}{}
besitzt der Punkt im Unendlichen die
\definitionsverweis {Multiplizität}{}{}


\mathl{d-1}{.}}

\faktzusatz {}


}

{


Die Gleichung für den projektiven Abschluss folgt direkt aus
Korollar 28.10.
Den Schnitt von $C$ mit der projektiven Geraden

\mathl{V_+(Z)}{}  im Unendlichen erhält man, wenn man in der Gleichung


\mavergleichskette

{\vergleichskette

{Z
}

{ = }{ 0
}

{  }{
}

{    }{
}

{   }{
}

}
{}{}{}
setzt. Bei


\mavergleichskette

{\vergleichskette

{d
}

{ = }{ 1
}

{  }{
}

{    }{
}

{   }{
}

}
{}{}{}
liegt insgesamt die Geradengleichung

\mathl{V_+(Y-aX-bZ)}{} vor, und der Schnitt mit

\mathl{V_+(Z)}{} legt den einzigen Punkt

\mathl{(1 , a ,0)}{} fest. Bei


\mavergleichskette

{\vergleichskette

{d
}

{ \geq }{ 2
}

{  }{
}

{    }{
}

{   }{
}

}
{}{}{}
liegt die Kurvengleichung


\mathdisp {V_+ \left(  YZ^{d-1} - s_dX^d -s_{d-1}X^{d-1}Z - \cdots -  s_0Z^d  \right)} {  }


mit


\mavergleichskette

{\vergleichskette

{s_d
}

{ \neq }{ 0
}

{  }{
}

{    }{
}

{   }{
}

}
{}{}{}
vor. Setzt man


\mavergleichskette

{\vergleichskette

{Z
}

{ = }{ 0
}

{  }{
}

{    }{
}

{   }{
}

}
{}{}{,}
so bleibt

\mathl{V_+ \left(  -s_dX^d  \right)}{} übrig, woraus


\mavergleichskette

{\vergleichskette

{X
}

{ = }{ 0
}

{  }{
}

{    }{
}

{   }{
}

}
{}{}{}
folgt. Dies entspricht dem einzigen unendlich fernen Punkt

\mathl{(0,1,0)}{.}


Für die Multiplizität betrachtet man die affine Gleichung der Kurve auf

\mathl{D_+(Y)}{.} D.h. man setzt


\mavergleichskette

{\vergleichskette

{ Y
}

{ = }{ 1
}

{  }{
}

{    }{
}

{   }{
}

}
{}{}{}
und erhält die affine Gleichung


\mathdisp {V \left(  Z^{d-1} - s_dX^d -s_{d-1}X^{d-1}Z - \cdots -  s_0Z^d  \right)} { , }


und der Punkt

\mathl{(0,1,0)}{} ist in diesen Koordinaten der Nullpunkt. Daher ist die Multiplizität gleich

\mathl{d-1}{} mit der einzigen durch


\mavergleichskette

{\vergleichskette

{ Z
}

{ = }{ 0
}

{  }{
}

{    }{
}

{   }{
}

}
{}{}{}
gegebenen Tangente. Bei


\mavergleichskette

{\vergleichskette

{ g
}

{ \geq }{ 3
}

{  }{
}

{    }{
}

{   }{
}

}
{}{}{}
ist die Multiplizität $\geq 2$ und daher liegt ein singulärer Punkt vor.


}


Dieser Satz ist so zu verstehen, dass bei


\mavergleichskette

{\vergleichskette

{d
}

{ \geq }{ 2
}

{  }{
}

{    }{
}

{   }{
}

}
{}{}{}
die $y$-Achse
\zusatzklammer {dafür steht der Punkt \mathlk{(0,1,0)}{}} {} {}
\anfuehrung{asymptotisch}{} zum Graphen gehört
\zusatzklammer {und auch die einzige Asymptote des Graphen ist} {} {.}
Die unendlich ferne Gerade

\mathl{V_+(Z)}{} ist die
\zusatzklammer {einzige} {} {}
Tangente an diesem Punkt. Die Normalisierung von $C$ ist der

\mathl{{\mathbb P}^{1}_{K}}{,} und zwar ist die Normalisierungsabbildung
nach Satz 29.5,
angewendet auf die affine Parametrisierung des Graphen
\maabbeledisp {} { {\mathbb A}^{1}_{K}} { {\mathbb A}^{1}_{K} \times {\mathbb A}^{1}_{K} =  {\mathbb A}^{2}_{K} \cong D_+(Z) \subset {\mathbb P}^{2}_{K}
} { x} { (x,F(x)) = (x,F(x),1)
} {,}
durch
\maabbeledisp {} {{\mathbb P}^{1}_{K}} {C \subset {\mathbb P}^{2}_{K}} { (x,t) } { \left(  xt^{d-1}, \hat{ F  }(x,t), t^d  \right)
} {}
gegeben. Dabei geht der unendlich ferne Punkt

\mathl{(1,0)}{} auf


\mavergleichskette

{\vergleichskette

{ (0,s_d,0)
}

{ = }{ (0,1,0)
}

{  }{
}

{    }{
}

{   }{
}

}
{}{}{.}


\inputfaktbeweis

{Ebene projektive Kurve/Graph einer rationalen Funktion in einer Variable/Singularität im Unendlichen/Fakt}

{Satz}

{}

{


\faktsituation {}

\faktvoraussetzung {Es sei $K$ ein
\definitionsverweis {algebraisch abgeschlossener Körper}{}{}
und seien


\mavergleichskette

{\vergleichskette

{ G,H
}

{ \in }{ K[X]
}

{  }{
}

{    }{
}

{   }{
}

}
{}{}{}
\definitionsverweis {Polynome in einer Variablen}{}{}
vom Grad


\mavergleichskette

{\vergleichskette

{ d , e
}

{ \geq }{ 1
}

{  }{
}

{    }{
}

{   }{
}

}
{}{}{}
ohne gemeinsame Nullstelle. Es sei


\mavergleichskette

{\vergleichskette

{ H
}

{ \neq }{ 0
}

{  }{
}

{    }{
}

{   }{
}

}
{}{}{}
und sei


\mavergleichskette

{\vergleichskette

{F(X)
}

{ = }{ G(X)/H(X)
}

{  }{
}

{    }{
}

{   }{
}

}
{}{}{}
die zugehörige rationale Funktion. Es seien

\mathl{\hat{  G   } (X,Z)}{} und

\mathl{\hat{  H   } (X,Z)}{} die zugehörigen
\definitionsverweis {Homogenisierungen}{}{}}

\faktfolgerung {Dann wird der
\definitionsverweis {projektive Abschluss}{}{}
$C$ des
\definitionsverweis {Graphen}{}{}
von

\mathl{F(X)}{} bei


\mavergleichskette

{\vergleichskette

{d
}

{ > }{ e
}

{  }{
}

{    }{
}

{   }{
}

}
{}{}{}
durch


\mathdisp {V_+ \left(  \hat{  H   } (X,Z)YZ^{d-e-1}- \hat{  G   } (X,Z)  \right)} {  }


und bei


\mavergleichskette

{\vergleichskette

{d
}

{ \leq }{e
}

{  }{
}

{    }{
}

{   }{
}

}
{}{}{}
durch


\mathdisp {V_+ \left(  \hat{  H   } (X,Z)Y- \hat{  G   } (X,Z)Z^{e-d+1}  \right)} {  }


beschrieben.}

\faktzusatz {}

\faktzusatz {}


}

{


Die affine Beschreibung der Kurve ist

\mathl{V(YH-G)}{.} Nach
Korollar 28.10 wird der projektive Abschluss durch die Homogenisierung von

\mathl{YH-G}{} beschrieben. Für diese ist der maximale Grad von \mathkon  { YH  }  {  und  }  {  G   }{ } ausschlaggebend, der Summand mit kleinerem Grad muss durch eine geeignete Potenz von $Z$ \anfuehrung{aufgefüllt}{} werden. Dies ergibt die angegebenen Gleichungen.


}


  \zwischenueberschrift{Monomiale projektive Kurven}


Zu einer ebenen monomialen Kurve

\mathl{s \mapsto (s^{e},s^{d})=(x,y)}{} mit teilerfremden Exponenten

\mathl{e >d}{} gehört nach
Satz 29.5
die monomiale projektive Kurve


\mathdisp {(s,t) \longmapsto (s^e,s^dt^{e-d} ,t^e)} { . }


Auf der offenen Menge

\mathl{D_+(t)}{} ist dies die ursprüngliche Abbildung und auf

\mathl{D_+(s)}{} ist dies die affine Abbildung


\mathdisp {t \longmapsto (t^{e-d}, t^e)} { . }


\inputfaktbeweis

{Ebene projektive monomiale Kurve/Singularität/Gesamtmultiplizität/Fakt}

{Satz}

{}

{


\faktsituation {}

\faktvoraussetzung {Es seien


\mavergleichskette

{\vergleichskette

{e
}

{ > }{ d
}

{  }{
}

{    }{
}

{   }{
}

}
{}{}{}
\definitionsverweis {teilerfremd}{}{.}
Für die durch


\mathdisp {(s,t) \longmapsto \left(  s^e  , \,  s^dt^{e-d}  , \,  t^e                 \right)} {  }


gegebene ebene projektive
\definitionsverweis {monomiale Kurve}{}{}
$C$ vom Grad $e$ gelten folgende Aussagen.}

\faktfolgerung {\aufzaehlungvier{Die Kurve wird durch die homogene Gleichung


\mavergleichskette

{\vergleichskette

{Y^e
}

{ = }{ X^dZ^{e-d}
}

{  }{
}

{    }{
}

{   }{
}

}
{}{}{}
vom Grad $e$ beschrieben.
}{Die Kurve ist
\definitionsverweis {glatt}{}{}
für alle Punkte \mathkon  { \neq (0,0,1)  }  {  und  }  { \neq (1,0,0)   }{ .}
}{Die Kurve hat im Punkt

\mathl{(0,0,1)}{} die
\definitionsverweis {Multiplizität}{}{}
$d$ und im Punkt

\mathl{(1,0,0)}{} die Multiplizität

\mathl{e-d}{.}
}{Bei


\mavergleichskette

{\vergleichskette

{ e
}

{ \geq }{ 3
}

{  }{
}

{    }{
}

{   }{
}

}
{}{}{}
ist die Kurve nicht glatt.
}}

\faktzusatz {}

\faktzusatz {}


}

{


\aufzaehlungvier{Die affine Gleichung ist

\mathl{X^d-Y^{e}}{,} und nach
Korollar 28.10
wird der projektive Abschluss durch die Homogenisierung, also durch

\mathl{V \left(  X^dZ^{e-d} - Y^{e}  \right)}{} beschrieben.
}{Auf der affinen Kurve


\mavergleichskette

{\vergleichskette

{ V(X^{d}-Y^{e})
}

{ \subset }{ {\mathbb A}^{2}_{ K }
}

{ \subset }{ {\mathbb P}^{2}_{ K }
}

{    }{
}

{   }{
}

}
{}{}{}
ist nach
Satz 20.12
nur der Nullpunkt, der dem projektiven Punkt

\mathl{(0,0,1)}{} entspricht,
\zusatzklammer {eventuell} {} {}
nicht glatt. Die Punkte auf der Kurve außerhalb von

\mathl{D_+(Z)}{} erhält man, indem man in der Gleichung


\mavergleichskette

{\vergleichskette

{ Z
}

{ = }{ 0
}

{  }{
}

{    }{
}

{   }{
}

}
{}{}{}
setzt. Dies erzwingt


\mavergleichskette

{\vergleichskette

{ Y
}

{ = }{ 0
}

{  }{
}

{    }{
}

{   }{
}

}
{}{}{,}
sodass es lediglich noch den Punkt

\mathl{(1,0,0)}{} gibt.
}{Die Multiplizität in einem Punkt ist eine lokale Eigenschaft. Der Punkt

\mathl{(0,0,1)}{} entspricht dem Nullpunkt auf der affinen monomialen Kurve


\mathl{V \left(  X^{d}-Y^e  \right)}{,} deren Multiplizität im Nullpunkt nach
Korollar 23.8
gleich dem kleineren Exponenten, also gleich $d$ ist. Der Punkt

\mathl{(1,0,0)}{} liegt auf

\mathl{D_+(X)}{} und dort ist

\mathl{V \left(  Y^{e}-Z^{e-d}  \right)}{} die affine Gleichung. Die Multiplizität ist wieder der kleinere Exponent, also gleich

\mathl{e-d}{.}
}{Folgt aus (3).
}


}
